% Pursuing Stacks (cotes 134-2 to 134-8): archive cote 134-5, site page 59 (PDF page 60), typescript page 246.
% typescript. Draft compared: ps4.tex:2984-3036.
% AI-transcribed by Claude Opus 5.5 (claude-opus-5-5, adaptive thinking, effort high), Message Batches API, batch ps2 (msgbatch_01NVwr1Yuo8qJf5rykeGBXPK), 2026-11-12; not yet checked by eye.
\scanpage{134-5/59}
\mspage{246}

seems to diminish somewhat the \xout{ax} emphasis I had put formerly upon the
notion of a test functor and it's weak variant, and enhance accordingly
the notion\xout{s} of an aspheric functor (with respect to a given asphericity
structure) and the two related asphericity notions for i which spring
from it (namely the notions of locally and \inserted{of} totally aspheric functors),
which now seem to come out as the more relevant and the more general ones.

To be wholly happy, we still need the relevant reformulation, in
terms of asphericity structures, of the main result of section 65, name-%
ly theorem 1 (p.177) characterizing test functors with values in (Cat),
under the assumption that the objects i(a) be contractible. This theorem,
I recall, has been the main outcome of the ``grimding'' reflections taking
their start with the``observation'' on p.94 about ten days earlier. The
approach then followed, as well as the contractibility assumption \over{on}{for} the
objects i(a) made in the theorem, retrospectively look awkward - it is
clear that the relevant notions of asphericity structures, and of asphe-%
ric functors into these, had been lacking. The theorem stated last
(p. 242) looks indeed a lot more satisfactory than the former, except
however in one respect - namely that the asphericity condition on i ,
in theorem 1 of p.177, can be expressed \xout{ax} by asphericity over $e_{A}{}^{\wedge}$ \uline{of}
\uline{just} $i^{*}(\hand{\Delta_{1}})$ \uline{alone}, rather than having to take $i^{*}(x)$ for all x in some
fixed generating subcategory of $(M,M_{as})$ containing i(A). To recover such
kind of \inserted[typed]{minimal} criterion in a more general case than (Cat), we'll have to rela-%
te the notion of an asphericity structure with the earlier one of a
\uline{contractibility} \uline{structure} ; \xout{which} \inserted[typed]{th\over{\unsure{is}}{e}} \inserted{latter} in our present reflections has faded
somewhat into the background, while at an earlier stage the homotopy and
contractibility notions had invaded the picture to an extent as to over-%
shadow and nearly bring to oblivion the magic of the ``asphericity game''.

% NOTE: the lower part of the page is hidden by a sheet laid over it (torn edge); only the top halves of the letters of the next line are visible
\unsure{We'll come back upon contractibility structures in the next section}

\margin{59}
